% =============================================================================
% Referee-revised patch v2: §3 sec:jet-limit-local-blocks
%
% Target file:   rh/rh_rebuild.tex
% Target lines:  474 through 578 (the entire current §3 block)
%
% Action: replace the existing §3 with the content below.
%
% Referee decision on the submitted patch:
%   - Same-point jet limit: accepted after making row/column ordering explicit.
%   - Cross-block jet limit: accepted after making row/column ordering explicit.
%   - Same-point Gram positivity: not accepted as originally written.
%     The determinant dominance used there does not follow from the current
%     (P1)--(P2) in Definition~\ref{def:local-phase-chart}. This revised
%     patch replaces that argument by an algebraic criterion plus an explicit
%     retained-packet Gram-regularity input.
%
% Status changes recommended:
%   \statusmig      not-started -> migrated
%   \statusreferee  not-started -> in-progress
%       (Taylor algebra passed; the Gram lower bound remains an explicit input)
%   \statussympy    not-started -> verified
%       (user-reported verification: rh/sympy/sec-jet-limit-local-blocks.py)
%   \statussim      not-started -> partial
%       (user-reported numerical sanity checks; not used as proof)
% =============================================================================

% =============================================================================
\section{Jet-limit local blocks}
\label{sec:jet-limit-local-blocks}
\statuspaper{LIVE}\quad
\statusmig{migrated}\quad
\statusreferee{in-progress}\quad
\statussympy{verified}\quad
\statussim{partial}\quad
\statuslean{not-started}
\par\medskip

The phase kernel \(K_\Ph\) collapses, in the limit \(h\to0\), into two
explicit \(2\times2\) matrices in the Gram-normalized basis of
\eqref{eq:point-to-jet-transform}: a same-point block \(J(T)\) attached
to a single center and a cross-block \(N_{12}(T_1,T_2)\) attached to two
centers.  These matrices, together with the same-point Gram-regularity input
isolated below, are the local blocks used downstream.  The phase enters only
through \(q,q',q''\) and \(\Del=\Ph(T_1)-\Ph(T_2)\).

\subsection{Same-point jet-limit}
\label{subsec:same-point-jet-limit}

\begin{lemma}[Same-point jet-limit]
\label{lem:same-point-jet-limit}
Let
\[
A_h(T)=
\begin{pmatrix}
K_\Ph(T-h,T-h) & K_\Ph(T-h,T+h)\\
K_\Ph(T+h,T-h) & K_\Ph(T+h,T+h)
\end{pmatrix},
\]
with row and column order \((T-h,T+h)\).  Then
\[
P_h A_h(T)P_h^\top \longrightarrow J(T)
\qquad (h\to0),
\]
where
\begin{equation}
\label{eq:same-point-J}
J(T)=\frac1\pi
\begin{pmatrix}
2q(T) & \tfrac12 q'(T)\\[1ex]
\tfrac12 q'(T) & \tfrac1{12}\bigl(q''(T)+2q(T)^3\bigr)
\end{pmatrix}.
\end{equation}
\end{lemma}

\begin{proof}
All derivatives in this proof are evaluated at \(T\).  Taylor expansion gives
\[
\Ph(T\pm h)=\Ph(T)\pm hq+\frac{h^2}{2}q'
\pm\frac{h^3}{6}q''+\frac{h^4}{24}q'''+O(h^5),
\]
and hence
\begin{equation}
\label{eq:phi-diff-same-point}
\Ph(T-h)-\Ph(T+h)=-2hq-\frac{h^3}{3}q''+O(h^5).
\end{equation}
Expanding the sine in the off-diagonal entry gives
\begin{equation}
\label{eq:A12-same-point-expansion}
K_\Ph(T-h,T+h)
=\frac{\sin(\Ph(T-h)-\Ph(T+h))}{-2\pi h}
=\frac{q}{\pi}+\frac{h^2}{6\pi}(q''-4q^3)+O(h^4).
\end{equation}
The diagonal entries are
\[
K_\Ph(T\pm h,T\pm h)
=\frac{q(T\pm h)}{\pi}
=\frac1\pi\left(q\pm hq'+\frac{h^2}{2}q''
\pm\frac{h^3}{6}q'''+O(h^4)\right).
\]
For any symmetric \(2\times2\) matrix \(M\), conjugation by
\eqref{eq:point-to-jet-transform} gives
\begin{align*}
(P_hMP_h^\top)_{11}
&=\frac12(M_{11}+2M_{12}+M_{22}),\\
(P_hMP_h^\top)_{12}
&=\frac1{4h}(M_{22}-M_{11}),\\
(P_hMP_h^\top)_{22}
&=\frac1{8h^2}(M_{11}-2M_{12}+M_{22}).
\end{align*}
Substituting the three entry expansions yields
\begin{align*}
(P_hA_hP_h^\top)_{11}
&=\frac{2q}{\pi}+\frac{2h^2}{3\pi}(q''-q^3)+O(h^4),\\
(P_hA_hP_h^\top)_{12}
&=\frac{q'}{2\pi}+\frac{h^2q'''}{12\pi}+O(h^4),\\
(P_hA_hP_h^\top)_{22}
&=\frac{q''+2q^3}{12\pi}+O(h^2).
\end{align*}
The singular powers in the lower row of \(P_h\) are exactly cancelled by the
symmetric Taylor differences.  Taking \(h\to0\) gives
\eqref{eq:same-point-J}.
\end{proof}

\begin{remark}[Symbolic verification]
\label{rem:same-point-jet-limit-sympy}
The entry-by-entry verification is automated in
\texttt{rh/sympy/sec-jet-limit-local-blocks.py}; the script checks
\[
\lim_{h\to0}P_hA_h(T)P_h^\top-J(T)=0
\]
as a symbolic identity in the Gram-normalized convention
\eqref{eq:point-to-jet-transform}.
\end{remark}

\subsection{Cross-block jet-limit}
\label{subsec:cross-block-jet-limit}

\begin{lemma}[Cross-block jet-limit]
\label{lem:cross-block-jet-limit}
For \(T_1\ne T_2\), set
\[
s=T_1-T_2,\qquad \Del=\Ph(T_1)-\Ph(T_2),\qquad q_j=q(T_j).
\]
Let
\[
C_h(T_1,T_2)=
\begin{pmatrix}
K_\Ph(T_1-h,T_2-h) & K_\Ph(T_1-h,T_2+h)\\
K_\Ph(T_1+h,T_2-h) & K_\Ph(T_1+h,T_2+h)
\end{pmatrix},
\]
with row order \((T_1-h,T_1+h)\) and column order
\((T_2-h,T_2+h)\).  Then
\[
P_h C_h(T_1,T_2)P_h^\top \longrightarrow
\frac1\pi N_{12}(T_1,T_2)
\qquad (h\to0),
\]
where
\begin{equation}
\label{eq:cross-N12}
N_{12}=
\begin{pmatrix}
\dfrac{2\sin\Del}{s}
&
\dfrac{\sin\Del-q_2s\cos\Del}{s^2}
\\[2.4ex]
\dfrac{q_1s\cos\Del-\sin\Del}{s^2}
&
\dfrac{(q_1+q_2)s\cos\Del+(q_1q_2s^2-2)\sin\Del}{2s^3}
\end{pmatrix}.
\end{equation}
\end{lemma}

\begin{proof}
Since \(T_1\ne T_2\), the kernel is smooth near \((T_1,T_2)\).  Write
\(K=K_\Ph(T_1,T_2)\), and use subscripts for derivatives at this point:
\begin{equation}
\label{eq:K-bivariate-taylor}
K_\Ph(T_1+u_1,T_2+u_2)
=K+u_1K_x+u_2K_y+\frac12u_1^2K_{xx}+u_1u_2K_{xy}
 +\frac12u_2^2K_{yy}+O(|u|^3).
\end{equation}
For a general \(2\times2\) matrix \(M\), conjugation by
\eqref{eq:point-to-jet-transform} gives
\begin{align*}
(P_hMP_h^\top)_{11}
&=\frac12(M_{11}+M_{12}+M_{21}+M_{22}),\\
(P_hMP_h^\top)_{12}
&=\frac1{4h}(-M_{11}+M_{12}-M_{21}+M_{22}),\\
(P_hMP_h^\top)_{21}
&=\frac1{4h}(-M_{11}-M_{12}+M_{21}+M_{22}),\\
(P_hMP_h^\top)_{22}
&=\frac1{8h^2}(M_{11}-M_{12}-M_{21}+M_{22}).
\end{align*}
Substitution of \eqref{eq:K-bivariate-taylor}, with
\(u_i\in\{-h,+h\}\), gives the derivative form
\begin{align*}
(P_hC_hP_h^\top)_{11}&\longrightarrow 2K,\\
(P_hC_hP_h^\top)_{12}&\longrightarrow K_y,\\
(P_hC_hP_h^\top)_{21}&\longrightarrow K_x,\\
(P_hC_hP_h^\top)_{22}&\longrightarrow \frac12K_{xy}.
\end{align*}
Direct differentiation of
\(K_\Ph(x,y)=\sin(\Ph(x)-\Ph(y))/\pi(x-y)\) at \((T_1,T_2)\) gives
\begin{align*}
K_\Ph(T_1,T_2)&=\frac{\sin\Del}{\pi s},\\
K_x(T_1,T_2)&=\frac{q_1s\cos\Del-\sin\Del}{\pi s^2},\\
K_y(T_1,T_2)&=\frac{\sin\Del-q_2s\cos\Del}{\pi s^2},\\
K_{xy}(T_1,T_2)
&=\frac{(q_1+q_2)s\cos\Del+(q_1q_2s^2-2)\sin\Del}{\pi s^3}.
\end{align*}
Combining these four identities gives the stated limit
\((1/\pi)N_{12}(T_1,T_2)\).
\end{proof}

\begin{remark}[Symbolic verification]
\label{rem:cross-block-jet-limit-sympy}
The entry-by-entry verification, including the bivariate expansion and the
explicit kernel-derivative evaluation, is automated in
\texttt{rh/sympy/sec-jet-limit-local-blocks.py}; the script checks
\[
\lim_{h\to0}P_hC_h(T_1,T_2)P_h^\top-(1/\pi)N_{12}(T_1,T_2)=0
\]
as a symbolic identity.
\end{remark}

\subsection{Same-point Gram algebra and retained-packet regularity}
\label{subsec:same-point-gram-positivity}

The preceding two lemmas are formal Taylor identities.  Positivity of the
same-point block is a separate retained-packet assertion unless it is proved
from the specific zeta phase chart.  The following algebra isolates exactly
what must be supplied.

\begin{definition}[Same-point Gram determinant]
\label{def:same-point-gram-determinant}
For the block \(J(T)\) in \eqref{eq:same-point-J}, define
\begin{equation}
\label{eq:same-point-gram-determinant}
D_J(T):=4q(T)^4+2q(T)q''(T)-3q'(T)^2.
\end{equation}
\end{definition}

\begin{lemma}[Algebraic same-point Gram criterion]
\label{lem:algebraic-same-point-gram-criterion}
The trace and determinant of \(J(T)\) are
\begin{align}
\label{eq:J-trace}
\Tr J(T)&=\frac{2q(T)^3+24q(T)+q''(T)}{12\pi},\\
\label{eq:J-determinant}
\det J(T)&=\frac{D_J(T)}{12\pi^2}.
\end{align}
Consequently \(J(T)\) is positive definite if and only if
\begin{equation}
\label{eq:same-point-PD-criterion}
q(T)>0
\qquad\text{and}\qquad
D_J(T)>0.
\end{equation}
\end{lemma}

\begin{proof}
The trace and determinant formulas follow directly from
\eqref{eq:same-point-J}.  Since \(J(T)\) is real symmetric and its leading
principal minor is \(2q(T)/\pi\), Sylvester's criterion gives positive
definiteness exactly when \(q(T)>0\) and \(\det J(T)>0\), equivalently
\eqref{eq:same-point-PD-criterion}.
\end{proof}

\begin{hypothesis}[Same-point Gram regularity on retained packets]
\label{hyp:same-point-gram-regularity}
On every retained packet there are absolute exponents \(C_0,C_1,C_D\) such
that
\begin{equation}
\label{eq:same-point-gram-regularity-upper}
|q(T)|+|q'(T)|+|q''(T)|\le Q^{C_0},
\end{equation}
\begin{equation}
\label{eq:same-point-gram-regularity-q-lower}
q(T)\ge Q^{-C_1},
\end{equation}
and
\begin{equation}
\label{eq:same-point-gram-regularity-det-lower}
D_J(T)=4q(T)^4+2q(T)q''(T)-3q'(T)^2\ge Q^{-C_D}.
\end{equation}
\end{hypothesis}

\begin{wip}\label{rem:wip-same-point-gram-positivity}
Hypothesis~\ref{hyp:same-point-gram-regularity}, especially the polynomial
lower bound \eqref{eq:same-point-gram-regularity-det-lower}, is not a formal
consequence of the current (P1)--(P2) in
Definition~\ref{def:local-phase-chart}.  It must either be proved during the
migration of Section~\ref{sec:local-kernel-and-jet-normalization}, folded
explicitly into the retained-packet definition, or replaced by a zeta-specific
asymptotic estimate for the chosen phase.  The algebraic trace/determinant
identities above are unconditional and are the portion verified symbolically.
\end{wip}

\begin{lemma}[Same-point Gram positivity]
\label{lem:same-point-gram-positivity}
Under Hypothesis~\ref{hyp:same-point-gram-regularity}, the block \(J(T)\) is
positive definite on retained packets and satisfies a polynomial lower
spectral bound
\[
\lambda_{\min}(J(T))\ge Q^{-C_J}
\]
for an absolute exponent \(C_J\).
\end{lemma}

\begin{proof}
Positive definiteness follows from \eqref{eq:same-point-PD-criterion}, using
\eqref{eq:same-point-gram-regularity-q-lower} and
\eqref{eq:same-point-gram-regularity-det-lower}.  Moreover
\[
\det J(T)=\frac{D_J(T)}{12\pi^2}\ge cQ^{-C_D}
\]
for an absolute constant \(c>0\).  The entrywise upper bound
\eqref{eq:same-point-gram-regularity-upper} and the explicit formula
\eqref{eq:same-point-J} give \(\lambda_{\max}(J(T))\le Q^{C'}\) for an
absolute exponent \(C'\).  Since
\(\lambda_{\min}(J(T))\lambda_{\max}(J(T))=\det J(T)\),
\[
\lambda_{\min}(J(T))
=\frac{\det J(T)}{\lambda_{\max}(J(T))}
\ge Q^{-C_J}
\]
with \(C_J=C_D+C'+O(1)\).
\end{proof}

\begin{remark}[Numerical check]
\label{rem:same-point-gram-positivity-numerical}
The numerical simulation script
\texttt{rh/simulation/sec-jet-limit-local-blocks.py} evaluates the
Riemann--Siegel theta phase at selected heights, checks positive definiteness
of \(J(T)\), and tests the convergence
\(P_hA_hP_h^\top\to J(T)\) and
\(P_hC_hP_h^\top\to(1/\pi)N_{12}\).  These computations are useful sanity
checks for the normalization and signs, but they are not used as a proof of
Hypothesis~\ref{hyp:same-point-gram-regularity}.
\end{remark}
